\BeleskeID{03-s020b}
% element: 03-s020b:e04
% element: 03-s020b:d02

Neka je \(H=I_3+I_2+I_1+I_0\) lokalno prisustvo zahteva. Pošto je u prikazanoj realizaciji \(GS=EI\,H\), izraz za prenos dozvole može se napisati kao
\[
 EO=EI\,\overline{GS}=EI\,\overline H.
\]
Poslednja jednakost sledi razmatranjem EI: kada je nula, EO je nula nezavisno od H; kada je jedinica, GS je jednak H. Zato dva stanja lokalnog zahteva pri zabrani daju isti EO, jednak nuli.

Na primer, neka viša grupa ima zahtev, srednja nema, a niža ima. Srednja dobija \(EI=0\). Inverzija njenog uslovljenog GS bila bi jedan i pogrešno bi dozvolila nižu grupu. Faktor EI u EO zadržava zabranu i kroz tu praznu srednju grupu.

Ako GS označava neuslovljeni lokalni zahtev H, drugi nivo kodera prioriteta i dalje bira najveći indeks među aktivnim grupama. Za formiranje nižih adresnih bitova, međutim, grupa bez lokalnog zahteva mora sačuvati informaciju da joj je viša grupa već zabranila rad. Zato je i u toj varijanti za prenos dozvole potreban izraz \(EI\,\overline H\).
